% ==============================================================================
% AUTOMATICALLY GENERATED BY SPECTRALBL REPORT PIPELINE
% GENERATED ON: 2026-06-12T18:41:40.967
% ACTIVE COMPUTATIONAL MAP OBJECT: TanhMap (Riemannian Vertical Manifold)
% CRITICAL: DO NOT MANUALLY AMEND THIS FILE. ALL CHANGES WILL BE OVERWRITTEN.
% ==============================================================================

\subsection{Analytical Coordinate Mapping and Riemannian Metric Formulations}
To project non-uniformly spaced atmospheric tower instrumentation from CASES-99 into a stable workspace for high-order pseudospectral methods, the pipeline instantiates a curvilinear coordinate chart modeled as a one-dimensional Riemannian manifold.

The physical coordinate $z \in [z_{\min}, z_{\max}]$ is linked to the computational coordinate $\xi \in [-1, 1]$ through smooth invertible maps $z = \mathcal{G}(\xi)$ and $\xi = \mathcal{F}(z)$.

\subsection{Symmetric Hyperbolic Tangent Transformation (\texttt{TanhMap})}
For centered shear-layer transitions, a symmetric hyperbolic tangent map with $\alpha = 2.5$ is applied.
The geometric midpoint is $z_c = 28.25\,\mathrm{m}$ and the span is $L = 53.5\,\mathrm{m}$.

\begin{equation}
    \xi = \mathcal{F}(z) = \frac{1}{\alpha} \operatorname{atanh}\left( \frac{2(z - z_c)\tanh(\alpha)}{L} \right)
\end{equation}
\begin{equation}
    z = \mathcal{G}(\xi) = z_c + \frac{L}{2} \frac{\tanh(\alpha \xi)}{\tanh(\alpha)}
\end{equation}
\begin{equation}
    \frac{dz}{d\xi} = \frac{L}{2} \left( \frac{\alpha}{\tanh(\alpha)} \right) \operatorname{sech}^2(\alpha \xi)
\end{equation}
\begin{equation}
    \frac{d^2z}{d\xi^2} = -L \left( \frac{\alpha^2}{\tanh(\alpha)} \right) \operatorname{sech}^2(\alpha \xi)\tanh(\alpha \xi)
\end{equation}


\subsubsection{Riemannian Metric Tensor and Volume Form}
The covariant and contravariant metric components are
\begin{equation}
    g_{\xi\xi} = J(\xi)^2 = \left(\frac{dz}{d\xi}\right)^2, \qquad
    g^{\xi\xi} = J(\xi)^{-2} = \left(\frac{dz}{d\xi}\right)^{-2}
\end{equation}
where $J(\xi)$ is the local Jacobian metric coefficient. The line element is $ds^2 = g_{\xi\xi} d\xi^2$.

Spatial variations in grid density are governed by the metric-gradient term
\begin{equation}
    J_{\xi}(\xi) = \frac{dJ}{d\xi} = \frac{d^2z}{d\xi^2}
\end{equation}
Because a one-dimensional manifold is intrinsically flat, $J_{\xi}(\xi)$ is a coordinate-stretching gradient, not intrinsic curvature. The induced volume form is
\begin{equation}
    dV = \sqrt{|g|} d\xi = J(\xi) d\xi
\end{equation}

\subsubsection{Geometric Inner Product and Hilbert Space Workspace}
The flow fields are represented in weighted $L^2_g([-1,1])$. For square-integrable fields $f$ and $g$,
\begin{equation}
    \langle f, g \rangle_g = \int_{-1}^{1} f(\xi) g(\xi) dV = \int_{-1}^{1} f(\xi) g(\xi) J(\xi) d\xi
\end{equation}
with induced norm $\|f\|_g^2 = \langle f, f \rangle_g$.

\subsubsection{Pseudospectral Operator Matrix Discretization}
Differentiation operators are assembled on the canonical Chebyshev computational domain using collocation nodes $\xi_j = \cos(\pi j / N)$ and then mapped to physical space:
\begin{equation}
    \frac{\partial \phi}{\partial z} = J^{-1} \frac{\partial \phi}{\partial \xi}
\end{equation}
\begin{equation}
    \frac{\partial^2 \phi}{\partial z^2} = J^{-2} \frac{\partial^2 \phi}{\partial \xi^2} - J^{-3} J_{\xi} \frac{\partial \phi}{\partial \xi}
\end{equation}

With $\mathbf{D}_{\xi} \in \mathbb{R}^{(N+1)\times(N+1)}$ and diagonal projections
\begin{equation}
    \mathbf{J}^{-1} = \operatorname{diag}\left(J^{-1}(\xi_j)\right), \quad
    \mathbf{J}^{-2} = \operatorname{diag}\left(J^{-2}(\xi_j)\right), \quad
    \mathbf{J}^{-3} = \operatorname{diag}\left(J^{-3}(\xi_j)\right)
\end{equation}
the discrete physical operators are
\begin{equation}
    \mathbf{D}_z = \mathbf{J}^{-1} \mathbf{D}_{\xi}
\end{equation}
\begin{equation}
    \mathbf{D}_{zz} = \mathbf{J}^{-2} \mathbf{D}_{\xi}^2 - \mathbf{J}^{-3} \mathbf{J}_{\xi} \mathbf{D}_{\xi}
\end{equation}
where $\mathbf{J}_{\xi} = \operatorname{diag}\left(J_{\xi}(\xi_j)\right)$.

\subsubsection{Geometric Quadrature Integration Consistency}
Continuous integrals transform as
\begin{equation}
    \int_{z_{\min}}^{z_{\max}} \phi(z) dz = \int_{-1}^{1} \phi\left(\mathcal{G}(\xi)\right) J(\xi) d\xi
\end{equation}
For Chebyshev-Gauss-Lobatto computational weights $w_j^{(\xi)}$, the physical weights are
\begin{equation}
    w_j^{(z)} = J(\xi_j) w_j^{(\xi)}
\end{equation}
This enforces dual consistency between metric-weighted differentiation and quadrature summations.
